\documentclass[11pt]{article}
\usepackage{amsmath,booktabs,multirow}
\usepackage{geometry}
\geometry{margin=1in}

\title{Supplementary Table S1: Representative Parameter Combinations Across Explored Parameter Space}
\author{}
\date{}

\begin{document}

\maketitle

\section*{Table S1: Parameter Sets and Corresponding Critical Thresholds}

This table documents representative parameter combinations spanning the biologically relevant parameter space explored in the main text. For each parameter set, we report the computed critical thresholds $\omega_{crit1}$ (generalist invasion) and $\omega_{crit2}$ (metabolite specialist displacement, when it exists).

\vspace{1em}

\begin{table}[h]
\centering
\small
\begin{tabular}{@{}lcccccccccc@{}}
\toprule
\multirow{2}{*}{Set} & \multicolumn{2}{c}{Growth rates} & \multicolumn{2}{c}{S-M mutualism} & \multicolumn{4}{c}{G cooperation} & \multirow{2}{*}{$\omega_{crit1}$} & \multirow{2}{*}{$\omega_{crit2}$} \\
\cmidrule(lr){2-3} \cmidrule(lr){4-5} \cmidrule(lr){6-9}
& $r_S$ & $r_M$ & $\sigma_{MS}$ & $\sigma_{SM}$ & $\sigma_{GS}$ & $\sigma_{GM}$ & $\sigma_{SG}$ & $\sigma_{MG}$ & & \\
\midrule
\textbf{Baseline} & 1.0 & 1.0 & 2.0 & 0.3 & 0.5 & 0.4 & 0.4 & 0.4 & 0.400 & -- \\
\textbf{Low mutualism} & 1.0 & 1.0 & 1.2 & 0.2 & 0.3 & 0.3 & 0.3 & 0.3 & 0.573 & -- \\
\textbf{High mutualism} & 1.0 & 1.0 & 2.8 & 0.3 & 0.6 & 0.5 & 0.5 & 0.5 & 0.287 & -- \\
\textbf{Strong G facilitation} & 1.0 & 1.0 & 2.0 & 0.3 & 0.7 & 0.6 & 0.5 & 0.8 & 0.325 & 0.686 \\
\textbf{Weak G-S cooperation} & 1.0 & 1.0 & 2.0 & 0.3 & 0.2 & 0.4 & 0.4 & 0.4 & 0.612 & -- \\
\textbf{Asymmetric rates} & 0.8 & 1.5 & 2.0 & 0.3 & 0.5 & 0.4 & 0.4 & 0.4 & 0.400 & -- \\
\textbf{High $\sigma_{MS}$} & 1.0 & 1.0 & 2.5 & 0.25 & 0.5 & 0.4 & 0.4 & 0.4 & 0.347 & -- \\
\textbf{Low $\sigma_{MS}$} & 1.0 & 1.0 & 1.5 & 0.3 & 0.5 & 0.4 & 0.4 & 0.4 & 0.478 & -- \\
\textbf{High $\sigma_{MG}$} & 1.0 & 1.0 & 2.0 & 0.3 & 0.5 & 0.5 & 0.4 & 1.2 & 0.380 & 0.612 \\
\textbf{Minimal window} & 1.0 & 1.0 & 2.0 & 0.3 & 0.5 & 0.5 & 0.4 & 1.0 & 0.371 & 0.531 \\
\midrule
\multicolumn{11}{l}{\textit{Competition parameters (constant across all sets unless noted):}} \\
\multicolumn{11}{l}{$\alpha_{GS} = 0.3$, $\alpha_{GM} = 0.3$, $\alpha_{SG} = 0.2$, $\alpha_{MG} = 0.2$} \\
\bottomrule
\end{tabular}
\end{table}

\section*{Notes:}

\begin{enumerate}
    \item \textbf{$\omega_{crit1}$}: Critical pathway allocation for generalist invasion into S-M equilibrium. Computed analytically using Eq. 3 in main text.

    \item \textbf{$\omega_{crit2}$}: Critical pathway allocation for metabolite specialist displacement from S-M-G equilibrium. Exists only when $\sigma_{MG}$ exceeds threshold (typically $\gtrsim 1.0$). Symbol "--" indicates $\omega_{crit2} \to \infty$ (permanent coexistence).

    \item \textbf{Parameter ranges explored systematically:}
    \begin{itemize}
        \item Growth rates: $r_S, r_M \in [0.5, 2.0]$ day$^{-1}$
        \item S-M mutualism: $\sigma_{MS} \in [1.1, 3.0]$, $\sigma_{SM} \in [0.2, 0.6]$
        \item G cooperation: $\sigma_{GS}, \sigma_{GM}, \sigma_{SG}, \sigma_{MG} \in [0.2, 1.2]$
        \item Competition: $\alpha_{GS}, \alpha_{GM}, \alpha_{SG}, \alpha_{MG} \in [0.1, 0.6]$
    \end{itemize}

    \item \textbf{Observed ranges across parameter space:}
    \begin{itemize}
        \item $\omega_{crit1} \in [0.15, 0.65]$ (minimum at high $\sigma_{MS}$ and $\sigma_{GS}$; maximum at low $\sigma_{MS}$ and high $\alpha_{GS}$)
        \item $\omega_{crit2} \in [0.50, 0.85]$ when it exists (minimum at low $\sigma_{MG}$; maximum at high $\sigma_{MG}$)
        \item Coexistence window width: $\Delta\omega = \omega_{crit2} - \omega_{crit1} \in [0.10, 0.35]$ when bounded
    \end{itemize}

    \item \textbf{Biological interpretation of parameter sets:}
    \begin{itemize}
        \item \textbf{Baseline}: Representative of moderately mutualistic cross-feeding systems (e.g., amino acid syntrophy in \textit{E. coli} co-cultures)
        \item \textbf{Low mutualism}: Minimal cross-feeding flux, approaching obligate dependence threshold ($\sigma_{MS} \approx 1$)
        \item \textbf{High mutualism}: Strong reciprocal benefits (e.g., tightly integrated co-cultures with engineered nutrient exchange)
        \item \textbf{Strong G facilitation}: Generalist provides substantial benefits to both specialists (large $\sigma_{MG}$), creating bounded coexistence window
        \item \textbf{Weak G-S cooperation}: Generalist-substrate cooperation limited, raising invasion threshold
        \item \textbf{Asymmetric rates}: Reflects natural systems with species having different doubling times
        \item \textbf{High/Low $\sigma_{MS}$}: Captures variation in substrate-to-metabolite cross-feeding flux
        \item \textbf{High $\sigma_{MG}$}: Generalist strongly facilitates metabolite specialist, creating second bifurcation
        \item \textbf{Minimal window}: Near the boundary of $\omega_{crit2}$ existence, $\Delta\omega \approx 0.16$
    \end{itemize}

    \item \textbf{Numerical verification:} All thresholds computed using analytical formulas (for $\omega_{crit1}$) and numerical root-finding (for $\omega_{crit2}$), validated via dynamical simulations showing bifurcation behavior.

    \item \textbf{Experimental mapping:} These parameter combinations can be realized experimentally by:
    \begin{itemize}
        \item Varying nutrient concentrations to modulate $\sigma$ parameters
        \item Using inducible promoters to control cross-feeding gene expression
        \item Selecting strains with different growth characteristics ($r$ values)
        \item Tuning inducer concentrations (e.g., arabinose/IPTG) to scan $\omega$ space
    \end{itemize}
\end{enumerate}

\section*{Parameter Sensitivity Analysis}

Gradient analysis across the explored parameter space reveals:

\begin{itemize}
    \item $\frac{\partial \omega_{crit1}}{\partial \sigma_{MS}} < 0$ (monotonic decrease): Stronger S-to-M mutualism consistently lowers invasion threshold
    \item $\frac{\partial \omega_{crit1}}{\partial \sigma_{GS}} < 0$: Increased S-to-G facilitation promotes generalist invasion
    \item $\frac{\partial \omega_{crit1}}{\partial \alpha_{GS}} > 0$: Competition with S raises invasion barrier
    \item $\frac{\partial \omega_{crit2}}{\partial \sigma_{MG}} < 0$ (when $\omega_{crit2}$ exists): G-to-M facilitation lowers displacement threshold, shrinking coexistence window
    \item $\frac{\partial \omega_{crit2}}{\partial \sigma_{MS}} > 0$: S-to-M mutualism raises displacement threshold, expanding coexistence window
\end{itemize}

These relationships are consistent across all tested parameter combinations, confirming robustness of the parameter landscape structure.

\section*{Validation Against Experimental Data}

Where available, parameters were constrained by experimental measurements:

\begin{itemize}
    \item Growth rates from batch culture doubling times \cite{harcombe2014metabolic,mee2014syntrophic}
    \item Mutualism strengths estimated from co-culture densities relative to monocultures \cite{shou2007synthetic,hoek2016resource}
    \item Pathway allocation $\omega$ mapped to inducer concentrations via dose-response curves \cite{guantes2006multistable}
\end{itemize}

The parameter ranges in this table encompass values reported in the experimental literature for engineered cross-feeding systems, validating biological relevance.

\end{document}
